% Pro C. Rabirio Postumo --- headnote
% pro-rabirio-postumo, 54 BC, Rome

\textit{For C. Rabirius Postumus, delivered at Rome in
54 or early 53 BC. Rabirius was a Roman financier --- the
son of C. Curtius, the foremost \textit{publicanus} of the
previous generation, adopted by the will of his maternal
uncle C. Rabirius, whose name he took. He had lent vast sums to the
exiled Egyptian king Ptolemy XII Auletes to finance the
king's restoration to his throne, a restoration carried
out in 55 BC by the proconsul of Syria, A. Gabinius. To
recover the debt Rabirius followed the king to Alexandria,
took the post of royal finance minister (\textit{dioecetes}),
and --- the points the prosecution most exploited --- wore
the Greek pallium and other un-Roman dress, and was for a
time held in the king's custody. The trial came in the
wake of Gabinius's own conviction for extortion under the
lex Iulia de repetundis: Gabinius had been condemned, the
damages assessed against him could not be recovered in
full, and Rabirius was now pursued under the clause of the
same law that allowed proceedings against those ``to whom
the money came'' (\textit{quo ea pecunia pervenerit}). The
prosecutor was C. Memmius; Cicero, who had reluctantly
defended Gabinius himself, now appeared for Rabirius.}

\textit{Cicero's case turns first on the legal mechanism.
The \textit{quo ea pecunia pervenerit} clause, he argues
(\textsection 8--13, 37), is no novelty but a transferred
chapter as old as the Servilian and Cornelian laws, and by
unbroken custom it had never reached anyone not first named
in the assessment of damages at the principal trial ---
whereas Rabirius's name appears nowhere in the Gabinian
record. Worse, the law on extortion does not bind the
equestrian order at all, and the heart of the speech
(\textsection 13--19) is an appeal to the Roman knights to
guard the immunity they hold from their fathers and not let
a new fire be set beneath their order. The defence of the
substance (\textsection 20--40) concedes the folly --- the
pallium, the foreign office, the lordship at Alexandria ---
but reframes it as the prior, fatal folly of having lent at
all: once entangled, Rabirius did under a king's compulsion
what Plato, Demetrius of Phalerum, and the consular P.
Rutilius Rufus had done before him. The speech closes
(\textsection 41--48) on C. Caesar, whose liberality alone
keeps the ruined Rabirius standing, and on Rabirius's old
kindness to Cicero in exile. No source records the verdict,
though the prominence of Caesar's support and the weakness
of the legal footing make an acquittal the likelier
outcome.}
